\chapter{Mathematische Grundlagen}

\section{Ziel dieses Bandes}

Die Projective Structure Theory (PST) erhebt den Anspruch, physikalische
Raumzeit nicht vorauszusetzen, sondern aus einer allgemeineren
informationellen Struktur entstehen zu lassen.

Während Band I die ontologischen Grundlagen des
Ontophänomenalismus (OP) behandelt, entwickelt dieser Band die
mathematische Sprache, welche für die Beschreibung des
Ontophänomenalen Phasenraums (OPP) benötigt wird.

Dabei verfolgen wir ein wichtiges Prinzip:

\begin{quote}
Es werden keinerlei geometrische Eigenschaften vorausgesetzt,
die später erst erklärt werden sollen.
\end{quote}

Insbesondere werden zunächst keine Größen wie

\[
\text{Raum},
\qquad
\text{Zeit},
\qquad
\text{Abstand},
\qquad
\text{Metrik}
\]

als fundamental angenommen.

Diese sollen vielmehr als mögliche Eigenschaften geeigneter
Projektionen aus dem OPP hervorgehen.

Damit unterscheidet sich die PST grundlegend von klassischen
geometrischen Theorien.

\section{Warum Mathematik?}

Die Physik benutzt Mathematik nicht nur als Beschreibungssprache,
sondern häufig als Werkzeug zur Entdeckung neuer Zusammenhänge.

Historisch entstanden zahlreiche mathematische Strukturen lange vor
ihrer physikalischen Anwendung.

Beispiele sind

\begin{itemize}
\item komplexe Zahlen,
\item Tensorrechnung,
\item Differentialgeometrie,
\item Lie-Gruppen,
\item Hilberträume,
\item Distributionen.
\end{itemize}

Die PST folgt derselben Idee.

Wir suchen keine Mathematik, welche bereits bekannte Physik beschreibt.

Wir suchen vielmehr mathematische Strukturen, aus denen bekannte
Physik emergieren kann.

\section{Minimale mathematische Anforderungen}

Da keine Geometrie vorausgesetzt werden soll, benötigt der OPP nur
wenige fundamentale Eigenschaften.

Es genügt zunächst eine Menge

\[
\Omega
\]

deren Elemente wir Attraktoren nennen.

Zwischen diesen Attraktoren existieren Beziehungen.

Mehr wird zunächst nicht angenommen.

Formal:

\[
\Omega
=
\{A_1,A_2,\ldots,A_n,\ldots\}.
\]

Es existiert keinerlei Ordnung.

Insbesondere existieren zunächst keine Begriffe wie

\[
A_i < A_j,
\]

keine Koordinaten

\[
(x,y,z),
\]

keine Zeitparameter

\[
t,
\]

und keine Abstände

\[
d(A_i,A_j).
\]

\section{Relationen}

Statt einer Geometrie wird ausschließlich angenommen, dass zwischen
zwei Attraktoren eine Relation existieren kann.

Formal schreiben wir

\[
R(A_i,A_j).
\]

Diese Relation besitzt zunächst keinerlei Interpretation.

Sie kann später beispielsweise darstellen:

\begin{itemize}
\item Ähnlichkeit,
\item Informationsaustausch,
\item Kopplungsstärke,
\item gegenseitige Beeinflussung,
\item Übergangswahrscheinlichkeit,
\item oder etwas vollständig anderes.
\end{itemize}

Der entscheidende Punkt lautet:

\begin{quote}
Die Relation ist fundamentaler als der Abstand.
\end{quote}

\section{Symmetrie der Relation}

Für die bisher untersuchten numerischen Modelle wurde stets eine
symmetrische Relation verwendet.

Es gilt daher

\[
R(A_i,A_j)
=
R(A_j,A_i).
\]

Dies bedeutet lediglich, dass die gegenseitige Beziehung unabhängig von
der Betrachtungsrichtung ist.

Eine Orientierung wird dadurch noch nicht eingeführt.

Asymmetrische Relationen können später untersucht werden.

Für die bisherigen Simulationen waren sie jedoch nicht notwendig.

\section{Selbstrelation}

Jeder Attraktor besitzt außerdem eine Selbstrelation.

Diese wird geschrieben als

\[
R(A_i,A_i).
\]

In sämtlichen bisherigen numerischen Experimenten wurde

\[
R(A_i,A_i)=1
\]

gewählt.

Dadurch erhält jeder Attraktor dieselbe maximale Selbstkopplung.

Dies entspricht keiner physikalischen Forderung, sondern lediglich einer
Normierung.

\section{Darstellung als Matrix}

Besitzt der OPP endlich viele Attraktoren,

\[
N
=
|\Omega|,
\]

dann kann die Relation vollständig durch eine Matrix beschrieben werden.

Es entsteht die Matrix

\[
\mathbf{R}
=
(R_{ij}),
\]

mit

\[
R_{ij}
=
R(A_i,A_j).
\]

Da

\[
R_{ij}=R_{ji},
\]

ist

\[
\mathbf{R}
=
\mathbf{R}^{\mathrm T}
\]

symmetrisch.

Diese Matrix bildet den Ausgangspunkt sämtlicher numerischer
Simulationen der PST.

Bemerkenswert ist dabei:

Die Matrix beschreibt ausschließlich Relationen.

Noch immer existieren keine Koordinaten.

Noch immer existiert kein Raum.

Noch immer existiert keine Zeit.

Die Matrix enthält ausschließlich informationelle Beziehungen.

\section{Vom Relationalen zur Geometrie}

Die zentrale Hoffnung der PST lautet nun:

Nicht die Matrix wird aus einer Geometrie berechnet.

Sondern umgekehrt.

Aus geeigneten Eigenschaften der Matrix

\[
\mathbf R
\]

sollen später emergieren:

\begin{align}
\mathbf R
&\longrightarrow
\text{Topologie}
\\
&\longrightarrow
\text{Nachbarschaft}
\\
&\longrightarrow
\text{Abstand}
\\
&\longrightarrow
\text{Metrik}
\\
&\longrightarrow
\text{Raumzeit}.
\end{align}

Dieser Perspektivwechsel bildet den eigentlichen Kern der Projective
Structure Theory.

\section{Bemerkung}

Die bisherigen numerischen Untersuchungen haben bereits gezeigt, dass
geeignete Relationen tatsächlich spektrale Strukturen erzeugen, welche
niedrige effektive Dimensionen besitzen können.

Damit existiert erstmals ein konkreter Hinweis darauf, dass
raumzeitähnliche Strukturen tatsächlich aus rein relationalen
Informationen hervorgehen könnten.

Die mathematische Beschreibung dieser spektralen Eigenschaften bildet den
Gegenstand des nächsten Kapitels.

